\documentclass[11pt]{article}
\usepackage{amsmath,amssymb,amsthm,hyperref}
\newtheorem{theorem}{Theorem}
\newtheorem{proposition}[theorem]{Proposition}
\newtheorem{lemma}[theorem]{Lemma}
\title{Auditing the Product Expansion in Selezneva's Construction}
\author{}
\date{24 November 2026}
\begin{document}
\maketitle

\paragraph{Status.}
The real product decomposition below is a complete written theorem, not a
new attention bound. The literature construction was inspected in full,
rather than inferred from its gate count. Its remaining transfer step
is identified explicitly. The original attention definitions are unchanged.

\paragraph{Sources inspected.}
S. N. Selezneva, ``Order of the length of Boolean functions in the class of
exclusive-OR sums of pseudoproducts'', Moscow University Computational
Mathematics and Cybernetics 40 (2016), 121--125,
\href{https://doi.org/10.3103/S0278641916030043}{doi:10.3103/S0278641916030043}.
The complete construction appears as Theorem 13.3, pages 154--155, in
Igor Sergeev, \emph{Etudes on the methods of designing efficient circuits},
\href{https://igorssergeev.github.io/papers/S-book_eng_280525.pdf}{author's book}.
The immediately preceding Lemma 13.2 and Corollary 13.1 give the
asymmetric hemisphere and quadrant covers used below. The circuit
operations in the original result are over the field with two elements.
In particular its input cube of monomial exponents must not be confused
with a partition of the actual input vertices.

\section*{A real version of the triangular elimination}
For positive integers $r,s$ let $x\in\{0,1\}^r$ and
$y\in\{0,1\}^s$. Identify a subset with its incidence vector and write
$x^A=\prod_{i\in A}x_i$. For centers $A\subseteq[r]$, $B\subseteq[s]$
define an exponent quadrant
\[
 Q(A,B)=\{(A\setminus\{i\},B\cup\{j\}):i\in A,\ j\notin B\}.
\]
Its covered exponents always belong to
$\mathcal E=\{(S,T):S\ne[r],\ T\ne\varnothing\}$.
Choose a family $\mathcal C$ of centers covering $\mathcal E$.
Order its centers by decreasing $|A|+|B|$, and break ties by
decreasing $|A|$. No particular ordering among the remaining ties is needed.

\begin{theorem}[Exact real product decomposition]
Every real function on the $(r+s)$-cube is the restriction of a sum of
at most
\[
 2|\mathcal C|+2^r+2^s
\]
scalar multiples of products of affine real forms. In each non-boundary
summand the factors have the more specific form
$y_j$ and $x_i+g_i(y)$, where $g_i$ is linear and uses only coordinates
outside the corresponding set $B$.
Consequently there is such a representation with $O(2^n/n^2)$ products
when $r=\lfloor n/2\rfloor$ and $s=n-r$.
\end{theorem}
\begin{proof}
Use the unique multilinear polynomial representing the given function.
First set aside its monomials outside $\mathcal E$. There are at most
$2^r+2^s$ such monomials, each already a product of affine forms.
Maintain a remainder supported on the part of $\mathcal E$ not yet
processed. At a center $(A,B)$ write its current coefficients on the
quadrant as $c_{ij}$, using zero for any exponent already removed, and set
\[
 g_i(y)=\sum_{j\notin B}c_{ij}y_j,
 \qquad
 E_{A,B}(x,y)=y^B\left(\prod_{i\in A}(x_i+g_i(y))-x^A\right).
 \tag{1}
\]
Expand the product. Terms choosing exactly one $g_i$ are precisely
\[
 \sum_{i\in A,\ j\notin B}c_{ij}x^{A\setminus\{i\}}y^{B\cup\{j\}},
 \tag{2}
\]
so subtracting (1) removes the desired quadrant coefficients. Every
remaining term chooses $k\ge2$ of the $g_i$. After replacing powers
$y_j^a$ with $y_j$ for $a\ge1$, its exponent $(S,T)$ satisfies
\[
 |S|=|A|-k\le |A|-2,\qquad
 |S|+|T|\le |A|+|B|,
 \qquad S\ne[r],\quad T\ne\varnothing.
 \tag{3}
\]
If this exponent belongs to another quadrant $Q(A',B')$, then
$|A'|=|S|+1$ and $|A'|+|B'|=|S|+|T|$. Thus its center occurs strictly
later in the stated order: either its total weight is smaller, or its
total weight is equal and $|A'|\le |A|-1$. No earlier removed coefficient
is changed. The maintained remainder condition follows by induction.
Once all covering centers have been processed the remainder is zero.
Each (1) is a difference of two products, proving the finite bound.

For the asymptotic assertion, the positive-hemisphere cover of the
$d$-cube excluding its zero vertex has $O(2^d/d)$ centers, as proved
in the cited Lemma 13.2. Complementation gives the negative-hemisphere
cover excluding the all-one vertex. Taking their Cartesian product
gives $|\mathcal C|=O(2^{r+s}/(rs))$. For balanced $r,s$, the two boundary
terms are $o(2^n/n^2)$, and $rs$ has order $n^2$.
\end{proof}

\section*{The actual cost of an exact ratio transfer}
The theorem packs an arbitrary matrix of real coefficients into a
product, which itself contains a linear number of independently
parameterized factors. It does not pack that matrix into one affine ratio.
This distinction has an exact parameter obstruction.

\begin{proposition}[Exact conversion of one product needs linearly many heads]
Let $r,s\ge1$ and put $n=r+s$. For matrices $c\in\mathbb R^{r\times s}$
consider
\[
 P_c(x,y)=\prod_{i=1}^r\left(x_i+\sum_{j=1}^s c_{ij}y_j\right)-x^{[r]}.
 \tag{4}
\]
If every $P_c$ has an exact representation on the Boolean cube by a
bias plus $H$ affine ratios with nonzero denominators, then
\[
 rs\le 1+2nH.
 \tag{5}
\]
This holds even without the attention denominator restrictions. In
particular a uniformly exact conversion of these balanced products
cannot use $o(n)$ heads.
\end{proposition}
\begin{proof}
The coefficient of $x^{[r]\setminus\{i\}}y_j$ in the unique multilinear
representative of (4) is exactly $c_{ij}$. Indeed the first-order terms
in the expansion give these coefficients, and terms selecting two or
more linear factors have strictly smaller $x$-degree.
Multilinear coefficients are fixed linear combinations of cube values.

Normalize each affine denominator to have constant term one. This is
permitted because the denominator is nonzero at the zero vertex.
Subtract each numerator's constant multiple of its denominator and
absorb it in the common bias. The score then has the form
\[
 b+\sum_{h=1}^H\frac{\sum_{i=1}^n u_{hi}z_i}
                            {1+\sum_{i=1}^n v_{hi}z_i},
 \tag{6}
\]
with $1+2nH$ real parameters. Evaluating (6) on cube vertices and
extracting the stated $rs$ coefficients gives a rational map in these
parameters. If $rs>1+2nH$, its coordinate rational functions are
algebraically dependent: after putting them over a common polynomial
denominator, compare the $\binom{rs+k}{rs}$ monomials of total degree
at most $k$ in the outputs with the $O(k^{1+2nH})$ possible monomials
in the parameters. For sufficiently large $k$ there is a nontrivial
linear dependence, giving a nonzero polynomial vanishing on the
rational image. Such a polynomial cannot vanish on all of
$\mathbb R^{rs}$. Exact conversion of every $P_c$ would make this map
surjective, a contradiction. This proves (5).
\end{proof}

\section*{What a sign transfer would still have to prove}
The preceding obstruction does not apply to signs. For example, if
$|c_{ij}|<1/2$, then at the actual cube points
$x=\boldsymbol1-e_i$, $y=e_j$, formula (4) gives
\[
 P_c(\boldsymbol1-e_i,e_j)
     =c_{ij}\prod_{k\ne i}(1+c_{kj}).
 \tag{7}
\]
All factors in the product on the right are positive. Arbitrary
$r\times s$ signs are therefore already present in one of the products
which would need to be compiled. This is only a test on a mixed
co-one-hot and one-hot grid. It is not a global construction.

Two separate unproved steps would be needed to turn (1) into the desired
attention bound. First, one would need a sublinear-cost conversion that
preserves appropriate signed values under the eventual sum, with every
denominator positive on the entire cube and all its nonzero slopes of
one strict sign. A sign test for each individual product does not
preserve signs of a sum of products. Second, any errors must respect
the triangular elimination despite the arbitrarily large and cancelling
real coefficients introduced by earlier steps. Neither follows from
the algebraic expansion. A theorem supplying both steps would itself
be a new universal compression mechanism. Without it, this route is
blocked at the conversion step, not close to an upper-bound proof.

\end{document}
